\documentclass[english,hebrew]{article}
\usepackage[HE8,T1]{fontenc}
\usepackage[utf8]{inputenc}
\setcounter{secnumdepth}{-2}
\setlength{\parindent}{0bp}
\usepackage{xcolor}
\usepackage{calc}
\usepackage{amsmath}
\usepackage{amsthm}
\usepackage{amssymb}

\makeatletter
\@ifundefined{date}{}{\date{}}
\usepackage{graphicx}
\usepackage{geometry}
\newcommand{\limi}{\underset{n\rightarrow\infty}{\lim}}
\newcommand{\limxz}{\underset{x\rightarrow x_{0}}{\lim}}
\newcommand{\limfxz}{\underset{x\rightarrow x_{0}}{\lim}f(x)}
\newcommand{\limfxm}{\underset{x\rightarrow x_{0}^{-}}{\lim}f(x)}
\newcommand{\limfxp}{\underset{x\rightarrow x_{0}^{+}}{\lim}f(x)}
\newcommand{\limxm}{\underset{x\rightarrow x_{0}^{-}}{\lim}}
\newcommand{\limxp}{\underset{x\rightarrow x_{0}^{+}}{\lim}}
\newcommand{\sqrm}{M_{m\times m}(\mathbb{F})}
\newcommand{\seqa}{(a_{n})_{n=1}^{\infty}}
\newcommand{\seqx}{(x_{n})_{n=1}^{\infty}}
\newcommand{\funcdr}{f:D\rightarrow\mathbb{R}}
\newcommand{\der}{\limxz\frac{f(x)-f(x_{0})}{x-x_{0}}}
\usepackage[all]{pdfprivacy}

\makeatother

\usepackage{babel}
\begin{document}
\newgeometry{top=1cm, bottom=2cm, left=1cm, right=1cm}
\title{אלגברה ליניארית {\beginL 2\endL} - הרצאה {\beginL 7\endL}}
\author{נעם לב}
\date{\L{$01.11.2026$}}

\maketitle
\noindent\fbox{\begin{minipage}[t]{1\columnwidth - 2\fboxsep - 2\fboxrule}%
\begin{center}
האוניברסיטה העברית \textbar{} אלגברה ליניארית ){\beginL 2\endL}( -
\L{$80135$} \textbar{} סמסטר{\small{} }{\beginL 2027\endL}{\small{} }א'
)תשפ\textquotedblleft ז א'( \textbar{} מרצה: פרופסור יבגני סטרחוב
\textbar{} מתרגל: מר לזר גוטרמן
\par\end{center}%
\end{minipage}}
\begin{center}
הסיכום הוא לא רשמי ולא עבר אישור של הסגל, על כן אין להשתמש בו )ובין
השאר במספור שלו ובשמות הטענות שבו( כמקור רשמי. אם ישנן טעויות, הן
שלי ולא של הסגל.
\par\end{center}

\medskip{}

\tableofcontents{}

\newpage{}

\section{תזכורת: הפולינום המינימלי וייצוג ע\textquotedblright י מטריצה משולשית
עליונה}

המשפט האחרון שהוכחנו )הרצאה \L{$6$}(:

יהי \L{$V$} מ\textquotedblright ו נ\textquotedblright ס מעל \L{$\mathbb{F}$}
כך ש- \L{%
\[
\dim V=n\ge1,\quad T\in\mathcal{L}(V)
\]
}%
אז התנאים הבאים שקולים:

\L{$(\text{I})$} קיים בסיס \L{$\mathcal{B}_{V}$} כך ש-\L{$M(T,\mathcal{B}_{V},\mathcal{B}_{V})$}
מטריצה משולשית עליונה.

\L{$(\text{II})$} \L{%
\[
q_{\min}(z)=(z-\alpha_{1})\cdots(z-\alpha_{m})
\]
}%
כאשר \L{$\alpha_{1},\dots,\alpha_{m}\in\mathbb{F}$}.

כלומר, אפשר להציג את הפולינום המינימלי כמכפלה של פולינומים מדרגה ראשונה.

דרשנו רק ש-\L{$\alpha_{1},\dots,\alpha_{m}$} יהיו מהשדה \L{$\mathbb{F}$},
ולא דרשנו שום דבר אחר )בפרט הם לא חייבים להיות שונים זה מזה(.

\section{שאלה: האם לסיבוב ב-\L{$\mathbb{R}^{2}$} יש ייצוג משולשי מעל \L{$\mathbb{R}$}?}

מ\textquotedblright ו ושדה:

\L{%
\[
\mathbb{F}=\mathbb{R},\quad V=\mathbb{R}^{2}
\]
}%

אופרטור:\L{%
\[
T((x,y))=(-y,x)
\]
}%

בסיס:\L{%
\[
\mathcal{B}_{V}=((1,0),(0,1))=(v_{1},v_{2})
\]
}%

מטריצה מייצגת של האופרטור לפי הבסיס:\L{%
\begin{align*}
T(v_{1}) & =T\left((1,0)\right)=\left(-0,1\right)=\left(0,1\right)=v_{2}=0_{\mathbb{F}}\cdot v_{1}+1\cdot v_{2}\\
T(v_{2}) & =T\left((0,1)\right)=\left(-1,0\right)=-v_{1}=-1\cdot v_{1}+0\cdot v_{2}
\end{align*}
}%
\L{%
\[
\therefore\quad M(T,\mathcal{B}_{V},\mathcal{B}_{V})=\begin{pmatrix}0 & -1\\
1 & 0
\end{pmatrix}
\]
}%

השאלה היא האם עבור אופרטור ומרחב אלה, שהם למעשה סיבוב של חצי פאי נגד
כיוון השעון סביב הנקודה \L{$\left(0,0\right)$}, קיים בסיס \L{$\widetilde{\mathcal{B}}_{V}$}
ב-\L{$V$} כך ש- \L{$M(T,\widetilde{\mathcal{B}}_{V},\widetilde{\mathcal{B}}_{V})$}
היא מטריצה משולשית עליונה.

שאלה זו שקולה לשאלה: האם קיימת מטריצה הפיכה \L{$Q$} כך ש- \L{%
\[
Q\begin{pmatrix}0 & -1\\
1 & 0
\end{pmatrix}Q^{-1}=\begin{pmatrix}* & *\\
0 & *
\end{pmatrix}
\]
}%

זו אותה שאלה, כי החלפת בסיס משנה את המטריצה המייצגת בדיוק כך, כאשר
\L{$Q$} היא מטריצת המעבר בין הבסיסים )לינארית \L{$1$}(.

\subsection{פתרון}

ראשית, נמצא את הפולינום המינימלי.

בהרצאה {\beginL 4\endL} מצאנו את הפולינום המינימלי של אופרטור זה מעל
\L{$\mathbb{C}$}.

אפשר לעבור על ההוכחה ולראות שבמקרה זה היא תקפה גם לאופרטור הזה מעל
\L{$\mathbb{R}$}.

אז\L{%
\[
q_{\min}(z)=z^{2}+1
\]
}%

אילו \L{$z^{2}+1=(z-\alpha_{1})(z-\alpha_{2})$} עם \L{$\alpha_{1},\alpha_{2}\in\mathbb{R}$},
אז הצבת \L{$z=\alpha_{1}$} נותנת \L{$\alpha_{1}^{2}=-1$}, וזאת סתירה.

לכן לא קיימים \L{$\alpha_{1},\alpha_{2}\in\mathbb{R}$} כך ש- \L{%
\[
q_{\min}(z)=(z-\alpha_{1})(z-\alpha_{2})
\]
}%
ולפי המשפט לא קיים בסיס \L{$\widetilde{\mathcal{B}}_{V}$} כך ש- \L{$M(T,\widetilde{\mathcal{B}}_{V},\widetilde{\mathcal{B}}_{V})$}
מטריצה משולשית עליונה.

השאלות האלו תמיד יהיו מעל השדה \L{$\mathbb{R}$} כי כפי שהוכחנו מעל
\L{$\mathbb{C}$} תמיד אפשר למצוא בסיס שבו המטריצה משולשית עליונה.

\section{דוגמה: אופרטור על \L{$\mathbb{R}^{2}$} עם פולינום אופייני \L{$(z-2)^{2}$}}

\L{%
\[
T\in\mathcal{L}(V),\quad V=\mathbb{R}^{2},\quad\mathbb{F}=\mathbb{R}
\]
}%
\L{%
\[
M(T,\mathcal{B}_{V},\mathcal{B}_{V})=\begin{pmatrix}1 & -1\\
1 & 3
\end{pmatrix}
\]
}%

\subsection{שאלה}

האם קיים בסיס \L{$\widetilde{\mathcal{B}}_{V}$} ב- \L{$V=\mathbb{R}^{2}$}
כך ש- \L{$M(T,\widetilde{\mathcal{B}}_{V},\widetilde{\mathcal{B}}_{V})$}
מטריצה משולשית עליונה?

\subsection{פתרון}

1. \L{%
\begin{align*}
P_{ch}(z) & =\det\begin{bmatrix}z-1 & 1\\
-1 & z-3
\end{bmatrix}\\
 & =(z-1)(z-3)+1\\
 & =z^{2}-4z+4\\
 & =(z-2)^{2}
\end{align*}
}%

2. 

מתקיים\L{%
\[
q_{\min}(z)=z-2
\]
}%
או \L{%
\begin{align*}
q_{\min}(z) & =(z-2)^{2}\\
 & =(z-2)(z-2)
\end{align*}
}%

כי הפולינום המינימלי מחלק את הפולינום האופייני ו-{\beginL 2\endL}
הוא ערך עצמי, לכן בהכרח {\beginL 2\endL} הוא שורש של הפולינום המינימלי
)לפי המשפטים שלמדנו בהרצאות חמש ושש(. אז אם הפולינום המינימלי ממעלה
{\beginL 1\endL} הוא חייב להיות \L{$z-2$}, ואם הוא ממעלה {\beginL 2\endL}
הוא חייב להיות הפולינום האופייני )ואם הוא מעל מעלה {\beginL 2\endL}
יש סתירה לכך שהוא מחלק את הפולינום האופייני(.

כלומר בכל מקרה אפשר להציג את הפולינום המינימלי כמכפלה של פולינומים
מדרגה ראשונה לכן קיים בסיס \L{$\widetilde{\mathcal{B}}_{V}$} כך ש-
\L{$M(T,\widetilde{\mathcal{B}}_{V},\widetilde{\mathcal{B}}_{V})$}
היא מטריצה משולשית עליונה וסיימנו.

\section{דוגמה: \L{$T(p)=2p+p'$} על \L{$\mathcal{P}_{3}(\mathbb{R})$} )\L{$\mathbb{R}_{\protect\leq3}[x]$}(.}

\L{%
\[
V=\mathcal{P}_{3}(\mathbb{R})
\]
}%
\L{%
\[
T(p(x))=2p(x)+p'(x)
\]
}%

\subsection{שאלה}

האם קיים בסיס \L{$\mathcal{B}_{V}$} במרחב \L{$V$} כך ש- \L{$M(T,\mathcal{B}_{V},\mathcal{B}_{V})$}
משולשית עליונה?\\

\textcolor{teal}{\L{$T$} הוא האגף השמאלי של המד\textquotedblright ר
\L{$p'+2p=q$}, ולכן הדוגמה מראה איך אפשר לתקוף מד\textquotedblright ר
לינאריות בכלים של אלגברה לינארית: במקום לפתור את המשוואה ישירות, חוקרים
את האופרטור \L{$T$} )ערכים עצמיים, פולינום מינימלי, ייצוג מטריציוני(.}

\subsection{פתרון}

1. \L{%
\[
T(p(x))=2p(x)+p'(x)
\]
}%
\L{%
\[
T(p(x))-2p(x)=p'(x)
\]
}%
\L{%
\[
\left(T-2I\right)\left(p(x)\right)=p'(x)
\]
}%

כלומר האופרטור \L{$T-2I$} שקול לפעולת הגזירה:\L{%
\[
T-2I=D=\frac{d}{dx}
\]
}%

הפולינום הוא ממעלה לכל היותר {\beginL 3\endL} ולכן גזירה שלו ארבע
פעמים תאפס אותו:\L{%
\[
\forall p\left(x\right)\in V:D^{4}(p)=0
\]
}%
\L{%
\[
\therefore\ (T-2I)^{4}=0
\]
}%

2. לעומת זאת, עבור למשל \L{$x^{3}\in V$} מתקיים\L{%
\[
D^{3}(x^{3})=6\neq0
\]
}%

)\L{$x^{3}\to3x^{2}\to6x\to6$}( כלומר לא לכל \L{$p\in\mathcal{P}_{3}(\mathbb{R})$}
מתקיים \L{$D^{3}(p)=0$}, דהיינו \L{$(T-2I)^{3}\neq0$} וזו לא העתקת
האפס.\\

3. קיבלנו שהפולינום \L{$(z-2)^{4}$} מאפס את \L{$T$}.

הפולינום המינימלי מחלק כל פולינום שמאפס את \L{$T$} )הרצאה \L{$5$}(,
ולכן \L{$q_{\min}(z)=(z-2)^{k}$} עבור \L{$k\in\left\{ 1,\dots,4\right\} $}
כלשהו.

נבדוק אותם.

אפשר להיווכח )כפי שעשינו ב-{\beginL 2\endL} עבור \L{$k=3$}( שעבור
\L{$k\in\left\{ 1,\dots,3\right\} $} הפולינום \L{$\left(z-2\right)^{k}$}
לא מאפס את האופרטור \L{$T$} ולכן נקבל\L{%
\begin{align*}
q_{\min}(z) & =(z-2)^{4}\\
 & =(z-2)(z-2)(z-2)(z-2)
\end{align*}
}%

4. לכן מהמשפט שלמדנו קיים בסיס \L{$\mathcal{B}_{V}$} שבו \L{$M(T,\mathcal{B}_{V},\mathcal{B}_{V})$}
משולשית עליונה.\\

\textcolor{teal}{לא תמיד אפשר לחשב את הפולינום המינימלי בקלות, אבל
כשיש שאלה על ייצוג ע\textquotedblright י מטריצה משולשית עליונה, כדאי
לנסות.}

\section{דוגמה: \L{$T=\lambda I$} לכסין עם ערך עצמי יחיד}

עכשיו אנחנו עוברים לשלב הבא, לא מתי יש ייצוג עם מטריצה משולשית עליונה
אלא מתי יש ייצוג עם מטריצה אלכסונית.

בהרצאה {\beginL 3\endL} למדנו שאם יש כ-\L{$\dim V$} ערכים עצמיים
שונים, אז אפשר לייצג את \L{$T$} ע\textquotedblright י מטריצה אלכסונית.

נשים לב שזה תנאי מספיק אבל לא הכרחי.

למשל, את האופרטור \L{$T=\lambda I$} על מ\textquotedblright ו ממימד
\L{$n$} )גם אם \L{$n\ge2$}( תמיד ניתן לייצג בעזרת מטריצה אלכסונית,
אפילו שיש רק ערך עצמי אחד.

\L{%
\[
M(T,\mathcal{B}_{V},\mathcal{B}_{V})=\begin{pmatrix}\lambda &  & 0_{\mathbb{F}}\\
 & \ddots\\
0_{\mathbb{F}} &  & \lambda
\end{pmatrix}
\]
}%

ולמעשה בכל בסיס המטריצה המייצגת של \L{$T$} אלכסונית.
\begin{proof}
יהי \L{$\mathcal{B}_{V}=(v_{1},\dots,v_{n})$} בסיס כלשהו.

לכל \L{$1\le j\le n$} מתקיים \L{%
\[
T(v_{j})=\lambda v_{j}=0_{\mathbb{F}}\cdot v_{1}+\dots+\lambda\cdot v_{j}+\dots+0_{\mathbb{F}}\cdot v_{n}
\]
}%
לכן בעמודה ה-\L{$j$} יש \L{$\lambda$} במקום ה-\L{$j$} ואפסים בשאר
המקומות, כלומר \L{%
\[
M(T,\mathcal{B}_{V},\mathcal{B}_{V})=\lambda I_{n}
\]
}%
שזו כמובן מטריצה אלכסונית.

\L{$\lambda$} הוא הערך העצמי היחיד: כל \L{$v\neq0_{V}$} מקיים \L{$T(v)=\lambda v$},
ואם \L{$T(v)=\mu v$} עם \L{$v\neq0_{V}$}, אז \L{$(\lambda-\mu)v=0_{V}$},
ולכן \L{$\mu=\lambda$}. 
\end{proof}

\section{משפט: תנאים שקולים ללכסינות}

יהי \L{$V$} מ\textquotedblright ו מעל שדה \L{$\mathbb{F}$} כלשהו
כך ש- \L{%
\[
\dim V=n\geq1,\quad T\in\mathcal{L}(V)
\]
}%
ויהיו \L{$\lambda_{1},\dots,\lambda_{m}$} כל הערכים העצמיים של \L{$T$}
)שונים זה מזה(, כלומר\L{%
\[
\forall i,j\in\{1,\dots,m\}\text{ s.t. }i\neq j:\quad\lambda_{j}\neq\lambda_{i}
\]
}%
אז התנאים הבאים שקולים:

\L{$(\text{I})$} קיים בסיס \L{$\mathcal{B}_{V}$} ב-\L{$V$} כך ש-
\L{%
\[
M(T,\mathcal{B}_{V},\mathcal{B}_{V})=\begin{pmatrix}a_{1,1} &  & 0_{\mathbb{F}}\\
 & \ddots\\
0_{\mathbb{F}} &  & a_{n,n}
\end{pmatrix}
\]
}%

\L{$(\text{II})$} קיים בסיס \L{$\mathcal{B}_{V}$} ב-\L{$V$} המורכב
מוקטורים עצמיים של \L{$T$}.

\L{$(\text{III})$} \L{%
\[
V=\text{Null}(T-\lambda_{1}I)\oplus\dots\oplus\text{Null}(T-\lambda_{m}I)
\]
}%

\L{$(\text{IV})$} \L{%
\[
\dim V=\dim\left(\text{Null}(T-\lambda_{1}I)\right)+\dots+\dim\left(\text{Null}(T-\lambda_{m}I)\right)
\]
}%

ברור כי \L{$m\leq n$}.

אבל זה ההבדל מהמשפט של הרצאה \L{$3$}, שם בהכרח \L{$m=n$} וכל הערכים
העצמיים שונים זה מזה.

כאן, על אף שהדרישה היא עדיין שהם שונים זה מזה, אין את הדרישה החזקה
על \L{$m$}: הוא יכול להיות אפילו \L{$1$}.

נזכיר ש{\beginL -\endL} \L{$\text{Null}(T-\lambda_{i}I)$} הוא המרחב
העצמי \L{$E_{\lambda_{i}}(T)$} מהרצאה \L{$1$}: כל הוקטורים העצמיים
השייכים ל- \L{$\lambda_{i}$} ו- \L{$0_{V}$}.

בהמשך נקבל משפט יותר פרקטי, דרך הפולינום האופייני.

נוכיח בצורה הבאה: \L{%
\[
(\text{I})\Rightarrow(\text{II})\Rightarrow(\text{III})\Rightarrow(\text{IV})\Rightarrow(\text{II})\Rightarrow(\text{I})
\]
}%

\subsection{הוכחה: \L{$(\text{I})\Rightarrow(\text{II})$}}

\L{$(\text{I})$} קיים בסיס \L{$\mathcal{B}_{V}$} ב-\L{$V$} כך ש-
\L{%
\[
M(T,\mathcal{B}_{V},\mathcal{B}_{V})=\begin{pmatrix}a_{1,1} &  & 0_{\mathbb{F}}\\
 & \ddots\\
0_{\mathbb{F}} &  & a_{n,n}
\end{pmatrix}
\]
}%

\L{$(\text{II})$} קיים בסיס \L{$\mathcal{B}_{V}$} ב-\L{$V$} המורכב
מוקטורים עצמיים של \L{$T$}.
\begin{proof}
1. יהי \L{$\mathcal{B}_{V}=(v_{1},\dots,v_{n})$} בסיס במרחב וקטורי
\L{$V$} כך ש- \L{%
\[
M(T,\mathcal{B}_{V},\mathcal{B}_{V})=\begin{pmatrix}a_{1,1} & 0_{\mathbb{F}} & \cdots & 0_{\mathbb{F}}\\
0_{\mathbb{F}} & a_{2,2} & \cdots & 0_{\mathbb{F}}\\
\vdots &  & \ddots & \vdots\\
0_{\mathbb{F}} & 0_{\mathbb{F}} & \cdots & a_{n,n}
\end{pmatrix}
\]
}%

2. אז\L{%
\begin{align*}
T(v_{1}) & =a_{1,1}v_{1}+0_{\mathbb{F}}v_{2}+\dots+0_{\mathbb{F}}v_{n}\\
T(v_{2}) & =0_{\mathbb{F}}\cdot v_{1}+a_{2,2}v_{2}+\dots+0_{\mathbb{F}}v_{n}\\
 & \ \vdots\\
T(v_{n}) & =0_{\mathbb{F}}\cdot v_{1}+0_{\mathbb{F}}\cdot v_{2}+\dots+a_{n,n}v_{n}
\end{align*}
}%
כלומר \L{%
\begin{align*}
T(v_{1}) & =a_{1,1}v_{1}\\
T(v_{2}) & =a_{2,2}v_{2}\\
 & \ \vdots\\
T(v_{n}) & =a_{n,n}v_{n}
\end{align*}
}%
בנוסף, \L{%
\[
v_{1}\neq0_{V},\dots,v_{n}\neq0_{V}
\]
}%
 כי \L{$\mathcal{B}_{V}=(v_{1},\dots,v_{n})$} הוא בסיס ולכן בפרט
בת\textquotedblright ל. 

אז \L{$v_{1},\dots,v_{n}$} הם וקטורים עצמיים של \L{$T$}.
\end{proof}

\subsection{הוכחה: \L{$(\text{II})\Rightarrow(\text{III})$}}

\L{$(\text{II})$} קיים בסיס \L{$\mathcal{B}_{V}$} ב-\L{$V$} המורכב
מוקטורים עצמיים של \L{$T$}.

\L{$(\text{III})$} \L{%
\[
V=\text{Null}(T-\lambda_{1}I)\oplus\dots\oplus\text{Null}(T-\lambda_{m}I)
\]
}%

\begin{proof}
1. יהי \L{$\mathcal{B}_{V}=(v_{1},\dots,v_{n})$} בסיס המורכב מוקטורים
עצמיים של \L{$T$}.

נסדר את הוקטורים מחדש לפי הערכים העצמיים שלהם.

לכל \L{$i\in\{1,\dots,m\}$}, נסמן את וקטורי הבסיס השייכים ל- \L{$\lambda_{i}$}
כך: \L{%
\[
\mathcal{B}^{(i)}=(v_{1}^{(i)},\dots,v_{n_{i}}^{(i)})
\]
}%
האינדקס העליון אומר לאיזה ערך עצמי הוקטור שייך, והתחתון ממספר את הוקטורים
בתוך הקבוצה. נשים את הקבוצות אחת אחרי השנייה: \L{%
\[
\widetilde{\mathcal{B}}_{V}=(\mathcal{B}^{(1)},\mathcal{B}^{(2)},\dots,\mathcal{B}^{(m)})
\]
}%
\L{$\widetilde{\mathcal{B}}_{V}$} הוא בסיס של \L{$V$}, כי אלה אותם
וקטורים בסדר אחר.

2. לא החסרנו ולא הוספנו וקטורים מהבסיס לכן\L{%
\[
n_{1}+\dots+n_{m}=n
\]
}%
אז לכל \L{$v\in V$}: \L{%
\[
v=\underbrace{\alpha_{1}^{(1)}v_{1}^{(1)}+\dots+\alpha_{n_{1}}^{(1)}v_{n_{1}}^{(1)}}_{u_{1}}+\dots+\underbrace{\alpha_{1}^{(m)}v_{1}^{(m)}+\dots+\alpha_{n_{m}}^{(m)}v_{n_{m}}^{(m)}}_{u_{m}}=u_{1}+\dots+u_{m}
\]
}%
כאשר \L{%
\[
u_{1}\in\text{Null}(T-\lambda_{1}I),\dots,u_{m}\in\text{Null}(T-\lambda_{m}I)
\]
}%
כי \L{$u_{i}$} הוא צ\textquotedblright ל של וקטורים מ- \L{$\text{Null}(T-\lambda_{i}I)$},
שהוא תמ\textquotedblright ו.\L{%
\[
\therefore\quad V=\text{Null}(T-\lambda_{1}I)+\dots+\text{Null}(T-\lambda_{m}I)
\]
}%

3. נרצה להוכיח שזה סכום ישר.

כדי לבדוק שסכום הוא ישר, מספיק להראות שהדרך היחידה להציג את \L{$0_{V}$}
כסכום כזה היא \L{$0_{V}+\dots+0_{V}$}.

נניח כי \L{%
\[
0_{V}=u_{1}+\dots+u_{m}
\]
}%
כאשר \L{%
\[
u_{1}\in\text{Null}(T-\lambda_{1}I),\dots,u_{m}\in\text{Null}(T-\lambda_{m}I)
\]
}%
ונניח בשלילה כי לא כל הוקטורים \L{$u_{1},\dots,u_{m}$} הם וקטור האפס.

הסדר לא משנה, ולכן אפשר למספר אותם מחדש כך שוקטורים שאינם וקטור האפס
יהיו הראשונים.

כלומר קיים \L{$p\geq1$} כך ש-\L{%
\[
\forall1\le i\le p:\quad u_{i}\ne0_{V}
\]
}%
\L{%
\[
\forall p+1\le j\le m:\quad u_{j}=0_{V}
\]
}%
נמחק מהסכום את האפסים: \L{%
\[
u_{1}+\dots+u_{p}=0_{V}
\]
}%
\L{$u_{1},\dots,u_{p}$} הם וקטורים עצמיים השייכים לערכים עצמיים שונים
)הם שייכים למרחבים שמשרים ערכים עצמיים שונים והם לא וקטור האפס(.

וקטורים עצמיים השייכים לערכים עצמיים שונים הם בת\textquotedblright ל
)הרצאה \L{$2$}(, ולכן הרשימה \L{$(u_{1},\dots,u_{p})$} בת\textquotedblright ל.

אבל \L{%
\[
1\cdot u_{1}+\dots+1\cdot u_{p}=0_{V}
\]
}%
וזאת סתירה לכך שהרשימה בת\textquotedblright ל.

לכן כל הוקטורים הם וקטור האפס.

כלומר, אם \L{%
\[
u_{1}+\dots+u_{m}=0_{V}\quad\wedge\quad u_{1}\in\text{Null}(T-\lambda_{1}I),\dots,u_{m}\in\text{Null}(T-\lambda_{m}I)
\]
}%
אז \L{%
\[
u_{1}=\dots=u_{m}=0_{V}
\]
}%
ולכן הסכום ישר: \L{%
\[
V=\text{Null}(T-\lambda_{1}I)\oplus\dots\oplus\text{Null}(T-\lambda_{m}I)
\]
}%
\end{proof}

\subsection{הוכחה: \L{$(\text{III})\Rightarrow(\text{IV})$}}

\L{$(\text{III})$} \L{%
\[
V=\text{Null}(T-\lambda_{1}I)\oplus\dots\oplus\text{Null}(T-\lambda_{m}I)
\]
}%

\L{$(\text{IV})$} \L{%
\[
\dim V=\dim\left(\text{Null}(T-\lambda_{1}I)\right)+\dots+\dim\left(\text{Null}(T-\lambda_{m}I)\right)
\]
}%

\begin{proof}
מלינארית {\beginL 1\endL}, המימד של סכום ישר שווה לסכום המימדים של
התמ\textquotedblright וים בסכום הישר, כנדרש.
\end{proof}

\subsection{הוכחה: \L{$(\text{IV})\Rightarrow(\text{II})$}}

\L{$(\text{IV})$} \L{%
\[
\dim V=\dim\left(\text{Null}(T-\lambda_{1}I)\right)+\dots+\dim\left(\text{Null}(T-\lambda_{m}I)\right)
\]
}%

\L{$(\text{II})$} קיים בסיס \L{$\mathcal{B}_{V}$} ב-\L{$V$} המורכב
מוקטורים עצמיים של \L{$T$}.
\begin{proof}
1. לכל \L{$i\in\{1,\dots,m\}$} נבחר בסיס של \L{$\text{Null}(T-\lambda_{i}I)$}
ונסמן אותו: \L{%
\[
\mathcal{B}^{(i)}=(v_{1}^{(i)},\dots,v_{n_{i}}^{(i)})
\]
}%
כלומר \L{$n_{i}=\dim\left(\text{Null}(T-\lambda_{i}I)\right)$}, וכל
הוקטורים ב-\L{$\mathcal{B}^{(i)}$} הם וקטורים עצמיים השייכים ל-\L{$\lambda_{i}$}
)הם שונים מ-\L{$0_{V}$} כי הם וקטורים בבסיס(.

ולכן לפי ההנחה \L{%
\[
n=n_{1}+\dots+n_{m}
\]
}%

2. נרצה להוכיח שהרשימה \L{%
\[
\mathcal{B}_{V}=(\mathcal{B}^{(1)},\dots,\mathcal{B}^{(m)})=(v_{1}^{(1)},\dots,v_{n_{1}}^{(1)},\dots,v_{1}^{(m)},\dots,v_{n_{m}}^{(m)})
\]
}%
היא בסיס של \L{$V$}.

3. אורך הרשימה הוא \L{$n_{1}+\dots+n_{m}=n=\dim V$}, ורשימה בת\textquotedblright ל
שאורכה \L{$\dim V$} היא בסיס של \L{$V$} )לינארית \L{$1$}(. 

לכן מספיק להוכיח שהרשימה \L{$\mathcal{B}_{V}$} בת\textquotedblright ל.

יהיו \L{$\alpha_{1}^{(1)},\dots,\alpha_{n_{1}}^{(1)},\dots,\dots,\alpha_{1}^{(m)},\dots,\alpha_{n_{m}}^{(m)}\in\mathbb{F}$}
סקלרים כך ש- \L{%
\[
\underset{u_{1}\in\text{Null}(T-\lambda_{1}I)}{\underbrace{\alpha_{1}^{(1)}v_{1}^{(1)}+\dots+\alpha_{n_{1}}^{(1)}v_{n_{1}}^{(1)}}}+\dots+\underset{u_{m}\in\text{Null}(T-\lambda_{m}I)}{\underbrace{\alpha_{1}^{(m)}v_{1}^{(m)}+\dots+\alpha_{n_{m}}^{(m)}v_{n_{m}}^{(m)}}}=0_{V}
\]
}%
ונראה שכל הסקלרים חייבים להיות \L{$0_{\mathbb{F}}$}.

נקבץ את האיברים לפי הקבוצות. 

לכל \L{$i\in\{1,\dots,m\}$} נסמן: \L{%
\[
u_{i}=\alpha_{1}^{(i)}v_{1}^{(i)}+\dots+\alpha_{n_{i}}^{(i)}v_{n_{i}}^{(i)}
\]
}%
אז \L{%
\[
u_{1}+\dots+u_{m}=0_{V}
\]
}%
ו- \L{$u_{i}\in\text{Null}(T-\lambda_{i}I)$}, כי הוא צירוף לינארי
של וקטורים מהתמ\textquotedblright ו הזה.

{\beginL 4\endL}. נראה ש- \L{%
\[
u_{1}=\dots=u_{m}=0_{V}
\]
}%
נניח בשלילה שלא כולם \L{$0_{V}$}.

הסדר לא משנה, ולכן אפשר למספר אותם מחדש כך שהוקטורים שאינם וקטור האפס
יהיו הראשונים.

כלומר קיים \L{$p\geq1$} כך ש-\L{%
\[
\forall1\le i\le p:\quad u_{i}\ne0_{V}
\]
}%
\L{%
\[
\forall p+1\le j\le m:\quad u_{j}=0_{V}
\]
}%
נמחק מהסכום את האפסים: \L{%
\[
u_{1}+\dots+u_{p}=0_{V}
\]
}%

\L{$u_{1},\dots,u_{p}$} הם וקטורים עצמיים השייכים לערכים עצמיים שונים
)הם שייכים למרחבים שמשרים ערכים עצמיים שונים והם לא וקטור האפס(.

וקטורים עצמיים השייכים לערכים עצמיים שונים הם בת\textquotedblright ל
)הרצאה \L{$2$}(, ולכן הרשימה \L{$(u_{1},\dots,u_{p})$} בת\textquotedblright ל.

אבל \L{%
\[
1\cdot u_{1}+\dots+1\cdot u_{p}=0_{V}
\]
}%
וזאת סתירה לכך שהרשימה בת\textquotedblright ל.

לכן כל הוקטורים הם וקטור האפס.

{\beginL 5\endL}. כלומר לכל \L{$i\in\{1,\dots,m\}$}: \L{%
\[
\alpha_{1}^{(i)}v_{1}^{(i)}+\dots+\alpha_{n_{i}}^{(i)}v_{n_{i}}^{(i)}=u_{i}=0_{V}
\]
}%
הרשימה \L{$\mathcal{B}^{(i)}$} היא בסיס של \L{$\text{Null}(T-\lambda_{i}I)$},
ובפרט בת\textquotedblright ל, ולכן: \L{%
\[
\alpha_{1}^{(i)}=\dots=\alpha_{n_{i}}^{(i)}=0_{\mathbb{F}}
\]
}%
כלומר כל המקדמים הם \L{$0_{\mathbb{F}}$}, והרשימה \L{$\mathcal{B}_{V}$}
בת\textquotedblright ל.

לסיכום, \L{$\mathcal{B}_{V}$} בסיס של \L{$V$}.

כל וקטור בו שייך לאחד הבסיסים \L{$\mathcal{B}^{(i)}$}, ולכן הוא וקטור
עצמי של \L{$T$} השייך ל-\L{$\lambda_{i}$}. 

כלומר \L{$(\text{II})$} מתקיים.
\end{proof}

\subsection{הוכחה: \L{$(\text{II})\Rightarrow(\text{I})$}}

\L{$(\text{II})$} קיים בסיס \L{$\mathcal{B}_{V}$} ב-\L{$V$} המורכב
מוקטורים עצמיים של \L{$T$}.

\L{$(\text{I})$} קיים בסיס \L{$\mathcal{B}_{V}$} ב-\L{$V$} כך ש-
\L{%
\[
M(T,\mathcal{B}_{V},\mathcal{B}_{V})=\begin{pmatrix}a_{1,1} &  & 0_{\mathbb{F}}\\
 & \ddots\\
0_{\mathbb{F}} &  & a_{n,n}
\end{pmatrix}
\]
}%

\begin{proof}
1. נניח כי \L{$\mathcal{B}_{V}=(v_{1},\dots,v_{n})$} בסיס של \L{$V$}
המורכב מוקטורים עצמיים של \L{$T$}.

כלומר, לכל \L{$i\in\{1,\dots,n\}$} קיים סקלר \L{$b_{i,i}\in\mathbb{F}$}
)הערך העצמי של \L{$v_{i}$}( כך ש- \L{%
\[
T(v_{i})=b_{i,i}v_{i}
\]
}%

2. נכתוב את זה כצירוף לינארי של כל וקטורי הבסיס: המקדם של \L{$v_{i}$}
הוא \L{$b_{i,i}$}, וכל שאר המקדמים הם \L{$0_{\mathbb{F}}$}: \L{%
\[
T(v_{i})=0_{\mathbb{F}}\cdot v_{1}+\dots+b_{i,i}\cdot v_{i}+\dots+0_{\mathbb{F}}\cdot v_{n}
\]
}%
העמודה ה-\L{$i$} של המטריצה המייצגת מורכבת מהמקדמים האלה: \L{$b_{i,i}$}
במקום ה-\L{$i$}, ו-\L{$0_{\mathbb{F}}$} בכל שאר המקומות. לכן: \L{%
\[
M(T,\mathcal{B}_{V},\mathcal{B}_{V})=\begin{pmatrix}b_{1,1} & 0_{\mathbb{F}} & \cdots & 0_{\mathbb{F}}\\
0_{\mathbb{F}} & b_{2,2} & \cdots & 0_{\mathbb{F}}\\
\vdots &  & \ddots & \vdots\\
0_{\mathbb{F}} & 0_{\mathbb{F}} & \cdots & b_{n,n}
\end{pmatrix}
\]
}%
כלומר המטריצה אלכסונית, ועל האלכסון מופיעים הערכים העצמיים. 
\end{proof}

\section{הערה - המשך הקורס}

מעל שדה כללי זה פחות או יותר סוף הסיפור.

מעל \L{$\mathbb{C}$} אפשר לומר עוד )צורת ז'ורדן(.

אם יהיה זמן בסוף הקורס נחזור לזה, אבל ההוכחה דורשת לפחות ארבע הרצאות
ואין לזה הרבה שימושים.

אולי במקום זה נדבר על ערכים סינגולריים, שחשובים יותר מבחינה פרקטית.

למשפט הספקטרלי נגיע בהחלט.

עוד מעט נתחיל לדבר על מרחבי מכפלה פנימית ועל אופרטורים במרחבי מכפלה
פנימית.
\end{document}
